% Variante independiente. Compilar desde la raíz: make el-bosque-docencia
\documentclass[11pt,letterpaper]{awesome-cv}
\usepackage{ifthen}
\usepackage[spanish]{babel}
\providecommand{\cvlang}{spanish}
\geometry{left=2cm,top=1.5cm,right=2cm,bottom=1.7cm,footskip=.5cm}
\fontdir[fonts/]
\colorlet{awesome}{awesome-red}
\renewcommand{\acvHeaderSocialSep}{\quad\textbar\quad}
\name{Oscar A.}{Quintero López}
\email{oquinterol@unal.edu.co}
\homepage{oquinterol.com}
\github{oquinterol}
\linkedin{oquinterol}
\orcid{0000-0001-8926-6400}
\position{Licenciado en Biología · Magíster en Bioinformática}
\makecvfooter{\today}{Oscar A. Quintero López · CV académico}{\thepage}
% Sin firma: incluirla solo por solicitud explícita del usuario.
% La clase original no inicia un párrafo tras el título de sección.
% Corrección local: evitar texto en la misma línea que el encabezado.
\renewenvironment{cvparagraph}{\par\vspace{2mm}\paragraphstyle}{\par\vspace{2mm}}

\begin{document}
\makecvheader

\cvsection{Perfil académico}
\begin{cvparagraph}
Licenciado en Biología y magíster en Bioinformática, con experiencia en análisis de datos genómicos y transcriptómicos, desarrollo de flujos de trabajo en Linux y reconstrucción metabólica. Trayectoria en apoyo a prácticas universitarias de Biología, capacitación de usuarios de laboratorio y acompañamiento de talleres de Linux y análisis de datos. Formación pedagógica y práctica profesional en educación en Biología, articuladas con herramientas computacionales para el estudio de sistemas biológicos.
\end{cvparagraph}

\ifthenelse{\equal{\cvlang}{spanish}}{
    \cvsection{Educación}
}{
    \cvsection{Education}
}

\begin{cventries}
    \ifthenelse{\equal{\cvlang}{spanish}}{
        \cventry
        {Magíster en Bioinformática} % Degree
        {Universidad Nacional de Colombia} % Institution
        {Bogotá D.C.---Colombia} % Location
        {Abr. 2026} % Fecha de grado según el expediente académico de septiembre de 2026
        {
            \begin{cvitems}
                \item {Tesis: Análisis genómico y reconstrucción metabólica de la variedad Colombia de \textit{Solanum tuberosum} Grupo Phureja.}
            \end{cvitems}
        }
    }{
        \cventry
        {M.Sc. in Bioinformatics} % Degree
        {National University of Colombia} % Institution
        {Bogotá D.C.---Colombia} % Location
        {Apr. 2026} % Graduation date in the September 2026 academic transcript
        {
            \begin{cvitems}
                \item {Thesis: Genomic Analysis and Metabolic Reconstruction of the Colombia Variety of \textit{Solanum tuberosum} Group Phureja.}
            \end{cvitems}
        }
    }

    \ifthenelse{\equal{\cvlang}{spanish}}{
        \cventry
        {Licenciado en Biología} % Degree
        {Universidad Distrital Francisco José de Caldas} % Institution
        {Bogotá D.C.---Colombia} % Location
        {Dic. 2022} % Fecha de grado según acta y diploma
        {
            \begin{cvitems}
                \item {Trabajo de grado: Desarrollo de flujo de trabajo (pipeline) para análisis de mutantes sólidos de \textit{Solanum tuberosum} Gr. Phureja.}
            \end{cvitems}
        }
    }{
        \cventry
        {Bachelor's Degree in Biology} % Degree
        {Francisco José de Caldas District University} % Institution
        {Bogotá D.C.---Colombia} % Location
        {Dec. 2022} % Graduation date in the degree certificate and diploma
        {
            \begin{cvitems}
                \item {Degree Project: Development of a Workflow (Pipeline) for the Analysis of Solid Mutants of \textit{Solanum tuberosum} Gr. Phureja.}
            \end{cvitems}
        }
    }

    \ifthenelse{\equal{\cvlang}{spanish}}{
        \cventry
        {Bachiller técnico con especialidad en Recreación y Deportes}
        {Institución Educativa Departamental Rufino Cuervo}
        {Chocontá, Cundinamarca---Colombia}
        {Dic. 2013}
        {}
    }{
        \cventry
        {Technical High School Diploma, Specialization in Recreation and Sports}
        {Rufino Cuervo Departmental Educational Institution}
        {Chocontá, Cundinamarca---Colombia}
        {Dec. 2013}
        {}
    }
\end{cventries}


\cvsection{Experiencia educativa y apoyo a la docencia}
\begin{cventries}
\cventry
{Técnico de Laboratorio Científico · Contratista}
{Universidad Distrital Francisco José de Caldas}
{Bogotá, Colombia}
{May. 2024---Sep. 2026}
{\begin{cvitems}
\item {Apoyo a prácticas de docencia e investigación en el Laboratorio de Biología de la Facultad de Ciencias Matemáticas y Naturales: alistamiento de equipos, insumos y espacios.}
\item {Capacitación de usuarios y orientación ante dificultades durante las prácticas; acompañamiento en el uso seguro de equipos e infraestructura.}
\item {Apoyo a la implementación, estandarización, validación y mejoramiento de técnicas, metodologías y protocolos para docencia, investigación y extensión.}
\end{cvitems}}

\cventry
{Asistente estudiantil · Especialización en Iluminación Pública y Privada}
{Universidad Nacional de Colombia}
{Bogotá, Colombia}
{Jun.---Dic. 2023}
{\begin{cvitems}
\item {Apoyo al desarrollo de laboratorios prácticos y a la coordinación de actividades académicas. Colaboración con el profesorado en la planificación y ejecución de actividades educativas y talleres.}
\end{cvitems}}

\cventry
{Práctica Profesional III y IV · Licenciatura en Biología}
{Instituto Técnico Industrial Francisco José de Caldas}
{Bogotá, Colombia}
{2020}
{\begin{cvitems}
\item {Práctica profesional en educación en Biología, jornada de la tarde, desarrollada en el marco de la Licenciatura de la Universidad Distrital.}
\end{cvitems}}
\end{cventries}

\newpage
\cvsection{Investigación en Bioinformática}
\begin{cventries}
\cventry
{Investigador en formación · Universidad Nacional de Colombia}
{GiBBS · Bioinformática y Biología de Sistemas}
{Bogotá, Colombia}
{Nov. 2023---Presente}
{\begin{cvitems}
\item {Desarrollo de un modelo metabólico a escala genómica de papa, integrando análisis genómico y biología de sistemas computacional.}
\item {Contribución a la redacción de propuestas de investigación y textos científicos relacionados con genómica y bioinformática.}
\end{cvitems}}

\cventry
{Investigador en formación · Universidad Distrital Francisco José de Caldas}
{BIOMOLc · Biología Molecular}
{Bogotá, Colombia}
{Ago. 2017---Presente}
{\begin{cvitems}
\item {Desarrollo de flujos de trabajo para análisis de RNA-Seq con herramientas basadas en Linux.}
\item {Apoyo a talleres de Linux básico y análisis de datos dentro del grupo de investigación.}
\end{cvitems}}
\end{cventries}

\cvsection{Competencias técnicas}
\begin{cvparagraph}
\textbf{Genómica y transcriptómica:} control de calidad con FastQC/MultiQC; alineamiento con BWA/Bowtie2; BLAST+; RNA-Seq; ensamblaje con hifiasm y Trinity; evaluación con BUSCO; anotación con AUGUSTUS; manejo de FASTA/FASTQ, BAM, VCF y GFF3.\par
\textbf{Biología de sistemas:} reconstrucción de modelos metabólicos a escala genómica; COBRApy, COBRA Toolbox, KEGG, SBML, MEMOTE y análisis FBA/FVA.\par
\textbf{Programación y análisis:} Python (pandas, NumPy, SciPy, Matplotlib), R (tidyverse/ggplot2), Bash y Jupyter.\par
\textbf{Trabajo reproducible:} Linux, Git, Conda, Docker, Snakemake, Nextflow y SLURM.
\end{cvparagraph}

\cvsection{Formación complementaria}
\begin{cventries}
\cventry
{Producción de proteínas recombinantes · Curso de 40 horas}
{Universidad Distrital · Comité Organizador CIIEB}
{Bogotá, Colombia}
{Nov. 2025}
{}
\cventry
{1.er Encuentro de Microbiología (EMIC) · Participación, 12 horas}
{American Society for Microbiology (ASM)}
{Cali, Colombia}
{Nov. 2025}
{}
\cventry
{Ensembl Browser Workshop}
{Ensembl}
{Bogotá, Colombia}
{Jul. 2024}
{}
\end{cventries}

\cvsection{Comunicación científica}
\begin{cvparagraph}
\textbf{Póster · ISCB-Latin America SoIBio CCBCOL 2024:} Ensamblaje de genoma de alta calidad de \textit{Solanum tuberosum} Grupo Phureja: análisis comparativo con el genoma público doble monoploide DM 1-3 516 R44.
\end{cvparagraph}
\begin{cvparagraph}
\textbf{Idiomas:} español nativo; inglés intermedio.
\end{cvparagraph}
\end{document}
